% !TeX program = pdflatex
%==============================================================================
%  PLANEJAMENTO DO CURSO DE APRENDIZADO DE MÁQUINA
%  Notas de aula -- UFF
%  Documento mestre: mapeia todas as aulas (oficiais e extras), seus tópicos
%  e as leituras correspondentes nos dois livros-texto adotados.
%==============================================================================
\documentclass[11pt,a4paper]{article}

\usepackage[utf8]{inputenc}
\usepackage[T1]{fontenc}
\usepackage{lmodern}
\usepackage{amsmath,amssymb,amsthm}
\usepackage[margin=2.5cm]{geometry}
\usepackage{longtable}
\usepackage{booktabs}
\usepackage{array}
\usepackage{enumitem}
\usepackage{xcolor}
\usepackage[colorlinks=true,linkcolor=blue,urlcolor=blue]{hyperref}

% Rótulos em português (babel-portugues indisponível nesta instalação)
\renewcommand{\abstractname}{Resumo}
\renewcommand{\contentsname}{Sumário}
\renewcommand{\refname}{Referências}
\renewcommand{\figurename}{Figura}
\renewcommand{\tablename}{Tabela}

\title{\textbf{Aprendizado de Máquina}\\[0.3em]
       \large Planejamento das Notas de Aula}
\author{Curso de Aprendizado de Máquina --- UFF\\
        \small Prof. Gabriel Sanfins\\
        \small Estatística, Ciências Atuariais, Matemática Aplicada e Engenharia Matemática}
\date{\today}

\begin{document}
\maketitle

\begin{abstract}
\noindent Este documento organiza todo o planejamento das notas de aula do curso.
Ele mapeia as \textbf{aulas oficiais} (acompanhando a sequência consolidada dos
slides do curso, numerados de 01 a 11) e as \textbf{aulas extras} (tópicos
complementares de aprendizado não supervisionado e aplicações), associando cada
aula aos respectivos capítulos dos dois livros-texto adotados. Para cada aula há
um arquivo \texttt{.tex} próprio nesta pasta, contendo o roteiro detalhado do que
deve ser dito em sala.
\end{abstract}

%==============================================================================
\section{Livros-texto adotados}
%==============================================================================
\begin{description}[leftmargin=2.2cm,style=nextline]
  \item[\textbf{[AME]}] Rafael Izbicki \& Tiago Mendonça dos Santos.
        \emph{Aprendizado de Máquina: Uma Abordagem Estatística}.
        Disponível on-line: \url{https://rafaelizbicki.com/ame/}.
  \item[\textbf{[ISLP]}] Gareth James, Daniela Witten, Trevor Hastie,
        Robert Tibshirani \& Jonathan Taylor.
        \emph{An Introduction to Statistical Learning, with Applications in Python}.
        Disponível on-line: \url{https://www.statlearning.com/}.
\end{description}

\noindent\textbf{Filosofia de uso dos dois livros.}
O [AME] é a espinha dorsal teórica: apresenta o aprendizado de máquina sob a ótica
estatística (função de risco, viés--variância, garantias teóricas), com notação
enxuta e tratamento unificado de regressão e classificação. O [ISLP] é a fonte de
intuição, exemplos e implementação em Python (\texttt{scikit-learn}), além de
fornecer figuras e conjuntos de dados clássicos. Nas notas, o esqueleto conceitual
segue o [AME] e os exemplos/laboratórios seguem o [ISLP].

%==============================================================================
\section{Estrutura geral do curso}
%==============================================================================
O curso está organizado em três blocos:
\begin{enumerate}[label=\textbf{\Roman*.}]
  \item \textbf{Fundamentos e Regressão} (Aulas 01--07): formalização do problema
        de aprendizado, regressão linear e regularização, seleção de modelos por
        validação cruzada, métodos não paramétricos (KNN, árvores, \emph{ensembles})
        e boas práticas de implementação (\emph{pipelines}).
  \item \textbf{Classificação} (Aulas 08--11): formulação do problema de
        classificação, classificadores gaussianos e logístico, métricas, SVM, KNN
        e árvores de classificação.
  \item \textbf{Aprendizado não supervisionado e aplicações} (Aulas Extras):
        agrupamento ($k$-médias), redução de dimensionalidade (PCA, t-SNE) e uma
        aplicação de processamento de linguagem natural (NLP) à classificação.
\end{enumerate}

\noindent Ao todo: \textbf{11 aulas oficiais} $+$ \textbf{3 aulas extras}
$=$ \textbf{14 conjuntos de notas de aula}, mais este documento de planejamento.

%==============================================================================
\section{Mapa das aulas}
%==============================================================================

\renewcommand{\arraystretch}{1.3}
\begin{longtable}{@{}p{0.6cm} p{4.6cm} p{5.0cm} p{4.2cm}@{}}
\toprule
\textbf{\#} & \textbf{Aula} & \textbf{Tópicos principais} & \textbf{Leitura} \\
\midrule
\endfirsthead
\toprule
\textbf{\#} & \textbf{Aula} & \textbf{Tópicos principais} & \textbf{Leitura} \\
\midrule
\endhead
\bottomrule
\endfoot

\multicolumn{4}{@{}l}{\textbf{Bloco I --- Fundamentos e Regressão}}\\[0.3em]

01 & Introdução ao Aprendizado de Máquina &
Notação e suposições; predição vs.\ inferência; as duas culturas; função de risco;
super e subajuste; balanço viés--variância; \emph{tuning parameters} &
AME cap.~1; ISLP cap.~1--2 \\

02 & Regressão Linear e Regularização &
Mínimos quadrados; além da linearidade; alta dimensão; seleção de subconjuntos e
\emph{stepwise}; Lasso e Ridge; esparsidade; interpretação bayesiana &
AME cap.~2--3; ISLP cap.~3, 6 \\

03 & Seleção de Modelos e Validação Cruzada &
Estimação do risco; \emph{data splitting}; validação cruzada com $k$ dobras;
escolha de $k$; \emph{bootstrap} &
AME §1.5; ISLP cap.~5 \\

04 & Métodos Não Paramétricos (KNN) &
Séries ortogonais e \emph{splines} (panorama); $k$ vizinhos mais próximos;
suavizadores lineares &
AME §4.1--4.3; ISLP §3.5, 7.1--7.5 \\

05 & Métodos Não Paramétricos: Aspectos Teóricos &
Maldição da dimensionalidade; taxas de convergência; esparsidade e redundância;
um pouco de teoria do KNN e séries ortogonais &
AME §4.13, cap.~5 \\

06 & Árvores de Regressão e \emph{Ensembles} &
Árvores de regressão (construção e poda); \emph{bagging}; florestas aleatórias;
importância de atributos; \emph{boosting} &
AME §4.8--4.10; ISLP cap.~8 \\

07 & Pré-processamento e \emph{Pipelines} &
Padronização (\texttt{StandardScaler}); vazamento de dados (\emph{data leakage});
construção de \emph{pipelines}; integração com validação cruzada &
ISLP labs cap.~5--6; AME §2.1.2 \\[0.6em]

\multicolumn{4}{@{}l}{\textbf{Bloco II --- Classificação}}\\[0.3em]

08 & Classificação e Classificadores Gaussianos &
Formulação e função de risco 0--1; classificadores \emph{plug-in}; regressão
logística; Bayes ingênuo; LDA e QDA &
AME cap.~7, §8.1; ISLP cap.~4 \\

09 & Métricas para Classificação &
Matriz de confusão; acurácia, precisão, \emph{recall}, $F_1$; curva ROC e AUC;
perdas assimétricas; dados desbalanceados; calibração &
AME §7.4, cap.~9; ISLP §4.4--4.5 \\

10 & Máquinas de Vetores de Suporte (SVM) &
Classificador de margem máxima; \emph{soft margin}; o truque do \emph{kernel};
SVM com kernels; relação com regressão logística; multiclasse &
AME §8.2, §4.6; ISLP cap.~9 \\

11 & KNN e Árvores de Classificação &
KNN para classificação; árvores de classificação (índice de Gini, entropia);
\emph{ensembles} para classificação; comparação entre classificadores &
AME §8.3, §8.6; ISLP §4.7.6, cap.~8 \\[0.6em]

\multicolumn{4}{@{}l}{\textbf{Bloco III --- Aulas Extras (Não supervisionado e Aplicações)}}\\[0.3em]

E1 & Análise de Agrupamento: $k$-Médias &
Aprendizado não supervisionado; $k$-médias (algoritmo de Lloyd); inicialização e
escolha de $k$; agrupamento hierárquico (panorama) &
AME cap.~11; ISLP §12.4 \\

E2 & Redução de Dimensionalidade (PCA e t-SNE) &
Componentes principais (PCA); interpretação geométrica; proporção de variância
explicada; aplicações (compressão); t-SNE para visualização &
AME cap.~10; ISLP §12.2 \\

E3 & NLP $+$ Classificação (Aplicação) &
Limpeza e exploração de texto; \emph{bag-of-words} / TF-IDF; classificação de
\emph{spam}; avaliação das métricas no contexto de texto &
Aplicação (notebook); AME §8.1.3 \\

\end{longtable}

%==============================================================================
\section{Aulas práticas e avaliações (contexto)}
%==============================================================================
As notas de aula são teóricas; o curso conta ainda com material de apoio ao qual
as notas fazem referência. Os \textbf{notebooks} ficam na pasta da própria aula,
em \texttt{aulas/}, e os materiais transversais em \texttt{recursos/}:
\begin{itemize}[leftmargin=1.4cm]
  \item \textbf{Aulas práticas} (notebooks): uma por aula,
        \texttt{Aula prática NN.ipynb}, na pasta da própria aula. São
        laboratórios guiados, para acompanhar em sala: texto e código
        alternados, uma ideia por célula, com blocos ``Sua vez'' para exercício.
        Cada uma reproduz as simulações das figuras da nota correspondente, com
        os mesmos parâmetros e a mesma semente.
  \item \textbf{Listas de exercícios}, também uma por aula e na pasta da própria
        aula, para depois da aula:
        \begin{itemize}
          \item \texttt{Lista teorica NN.tex/.pdf} --- 3 ou 4 exercícios de
                nível fácil a mediano, todos tratados igual.
                O gabarito sai do \emph{mesmo} arquivo-fonte
                (\texttt{Lista teorica NN - gabarito.pdf}), o que impede que
                enunciado e solução divirjam;
          \item \texttt{Lista prática NN.ipynb} --- notebook com lacunas
                marcadas por \texttt{...} para o aluno completar, acompanhado de
                \texttt{Lista prática NN - gabarito.ipynb}.
        \end{itemize}
  \item \textbf{Exemplos} em Python, na pasta da aula correspondente: comparação
        de classificadores paramétricos (aula 08), SVM (10), KNN e árvores (11),
        $k$-médias (E1) e PCA (E2).
  \item Duas \textbf{avaliações presenciais} (AP1 e AP2) com gabarito, em
        \texttt{recursos/avaliacoes}.
\end{itemize}

%==============================================================================
\section{Convenções das notas de aula}
%==============================================================================
\begin{itemize}[leftmargin=1.4cm]
  \item Cada aula tem um arquivo \texttt{NN Título.tex} próprio, compilável a
        partir da pasta da própria aula. O preâmbulo é compartilhado: cada
        \texttt{.tex} carrega o estilo comum com
        \verb|\usepackage{../../recursos/latex/estilo-notas}|, de modo que
        notação, ambientes de teorema e caixas pedagógicas ficam definidos em um
        único lugar.
  \item Estrutura padrão de cada nota: \emph{objetivos da aula} $\to$
        \emph{motivação} $\to$ \emph{desenvolvimento teórico} (definições,
        resultados, demonstrações quando cabível) $\to$ \emph{intuições e exemplos}
        $\to$ \emph{prática em Python} (referência ao notebook) $\to$
        \emph{resumo} e \emph{leitura recomendada}.
  \item Notação unificada seguindo o [AME]: $(X,Y)$ par aleatório,
        $r(x)=\mathbb{E}[Y\mid X=x]$ função de regressão, $g(x)$ classificador,
        $\mathcal{D}=\{(X_i,Y_i)\}_{i=1}^n$ amostra de treino, $R(\cdot)$ função de
        risco.
  \item Caixas de \emph{``o que dizer em sala''} destacam o fio condutor do discurso
        e as perguntas a fazer aos alunos.
\end{itemize}

\end{document}
